%%%%%%%%%%%%%%%%%%%%%%%%%%%%%%%%%%%%%%%%%%%%%%%%%%%%%%%%%%%%%%%%%%%%%%%%%%%%%%%
%%  TKS_v7.4_CoreCalculus.tex
%%  Core Calculus Overview: Minimal Central Engine for Academic Readers
%%
%%  Part of TKS v7.4 --- To be placed in Part 0 (Canonical Foundation)
%%  AFTER Chapter: Ontological Foundations Reference
%%%%%%%%%%%%%%%%%%%%%%%%%%%%%%%%%%%%%%%%%%%%%%%%%%%%%%%%%%%%%%%%%%%%%%%%%%%%%%%

\chapter{Core Calculus Overview}
\label{ch:core-calculus}

\begin{tcolorbox}[colback=cyan!5!white,colframe=cyan!75!black,title=\textbf{Purpose}]
This chapter provides a self-contained summary of the \emph{minimal} TKS formal system
for academics unfamiliar with the full manual. It distills the essential objects,
algebraic structures, and core theorems into a concise reference suitable for
researchers seeking rapid orientation before engaging with advanced material.
All definitions herein are consistent with and derived from the canonical TKS v6.1--v7.4
formalization.
\end{tcolorbox}

%% ============================================================================
%% SECTION 1: CORE OBJECTS
%% ============================================================================
\section{Core Objects}
\label{sec:core-objects}

The TKS formal system is built upon six fundamental object classes: Worlds, Elements,
Noetics, Foundations, Sub-Foundations, and Acquisitions. We define each precisely.

\begin{definition}[Worlds]
\label{def:core-worlds}
The \textbf{World set} $\Worlds$ is the finite totally ordered set:
\[
    \Worlds := \{A, B, C, D\}
\]
equipped with the ordering $A \succ B \succ C \succ D$, representing decreasing
ontological subtlety. Each world carries a distinct interpretation:
\begin{center}
\begin{tabular}{clll}
\toprule
\textbf{World} & \textbf{Label} & \textbf{Domain} & \textbf{Ordinal} \\
\midrule
$A$ & Archetypal & Spiritual / Divine & $\mathrm{ord}(A) = 0$ \\
$B$ & Intellectual & Mental / Conceptual & $\mathrm{ord}(B) = 1$ \\
$C$ & Conceptual & Emotional / Formative & $\mathrm{ord}(C) = 2$ \\
$D$ & Material & Physical / Manifest & $\mathrm{ord}(D) = 3$ \\
\bottomrule
\end{tabular}
\end{center}
The \emph{world distance} $d(W_1, W_2) := |\mathrm{ord}(W_1) - \mathrm{ord}(W_2)|$
quantifies separation between ontological strata.
\end{definition}

\begin{definition}[Elements]
\label{def:core-elements}
The \textbf{Element space} $\Elements$ is the Cartesian product:
\[
    \Elements := \Worlds \times \{1, 2, 3, 4, 5, 6, 7, 8, 9, 10\}
\]
yielding $|\Elements| = 40$ atomic constituents. An element $Wn$ (where $W \in \Worlds$
and $n \in \{1, \ldots, 10\}$) combines a world stratum with a \emph{mode} encoding
its functional role. The ten modes are:
\begin{center}
\begin{tabular}{cl|cl}
\toprule
\textbf{Mode} & \textbf{Role} & \textbf{Mode} & \textbf{Role} \\
\midrule
1 & Mind / Awareness & 6 & Male / Projective \\
2 & Positive / Attraction & 7 & Rhythm / Periodicity \\
3 & Negative / Repulsion & 8 & Above / Causal / Inner \\
4 & Vibration / Energy & 9 & Below / Effect / Outer \\
5 & Female / Receptive & 10 & Idea / Pure Potential \\
\bottomrule
\end{tabular}
\end{center}
\end{definition}

\begin{definition}[Noetics]
\label{def:core-noetics}
The \textbf{Noetic operator space} $\Noetics$ is the set of ten consciousness operators:
\[
    \Noetics := \{\noe{0}, \noe{1}, \noe{2}, \noe{3}, \noe{4}, \noe{5},
                  \noe{6}, \noe{7}, \noe{8}, \noe{9}\}
\]
Each noetic $\noe{k}$ is an endofunction on the Idea space:
\[
    \noe{k} : \Ideas \longrightarrow \Ideas
\]
The noetic $\noe{0}$ (Idea) acts as the identity. Noetics form three \emph{dual pairs}
summing to 9: $(\noe{2}, \noe{7})$, $(\noe{3}, \noe{6})$, $(\noe{4}, \noe{5})$,
plus the distinguished pairs $(\noe{1}, \noe{8})$ (Mind--Above). Composition of
dual pairs approximately yields the identity: $\noe{j} \circ \noe{i} \approx \noe{0}$
when $i + j = 9$.
\end{definition}

\begin{definition}[Foundations]
\label{def:core-foundations}
The \textbf{Foundation space} $\Foundations$ consists of seven organizational principles:
\[
    \Foundations := \{\Foundations_1, \Foundations_2, \Foundations_3, \Foundations_4,
                      \Foundations_5, \Foundations_6, \Foundations_7\}
\]
Each Foundation $\Foundations_m$ governs a life domain:
\begin{center}
\begin{tabular}{cll}
\toprule
\textbf{Index} & \textbf{Foundation} & \textbf{Formal Principle} \\
\midrule
$\Foundations_1$ & Unity / Association & \textsc{Coherence} \\
$\Foundations_2$ & Wisdom & \textsc{Epistemic-Structure} \\
$\Foundations_3$ & Life & \textsc{Vitality} \\
$\Foundations_4$ & Companionship & \textsc{Relational-Binding} \\
$\Foundations_5$ & Power & \textsc{Causal-Efficacy} \\
$\Foundations_6$ & Material / Wealth & \textsc{Resource-Structure} \\
$\Foundations_7$ & Lust / Continuation & \textsc{Generative-Drive} \\
\bottomrule
\end{tabular}
\end{center}
\end{definition}

\begin{definition}[Sub-Foundations]
\label{def:core-subfoundations}
The \textbf{Sub-Foundation space} is the Cartesian product:
\[
    \Foundations^* := \Foundations \times \{a, b, c, d\}
\]
yielding $|\Foundations^*| = 28$ refinements, where the second index indicates
world-level manifestation ($a$ = Spiritual, $b$ = Mental, $c$ = Emotional, $d$ = Physical).
For example, $\Foundations_{3d}$ denotes Physical Life (biological vitality).
\end{definition}

\begin{definition}[Acquisitions]
\label{def:core-acquisitions}
The \textbf{Acquisition space} $\Acquisitions$ consists of 22 derived composite structures:
\[
    \Acquisitions := \{A_0\} \cup \{D_m\}_{m=1}^{7} \cup \{W_m\}_{m=1}^{7}
                     \cup \{P_m\}_{m=1}^{7}
\]
where:
\begin{itemize}
    \item $A_0$: Pure Desire (root intentional vector),
    \item $D_m$: Desire for Foundation $\Foundations_m$,
    \item $W_m$: Wisdom about Foundation $\Foundations_m$,
    \item $P_m$: Power to achieve Foundation $\Foundations_m$.
\end{itemize}
An \emph{acquisition state} $\sigma : \Acquisitions \to \{\mathtt{true}, \mathtt{false}\}$
tracks which acquisitions are satisfied.
\end{definition}

%% ============================================================================
%% SECTION 2: ALGEBRAIC STRUCTURE
%% ============================================================================
\section{Algebraic Structure}
\label{sec:algebraic-structure}

We now specify the algebraic structures governing TKS: the monoid of Ideas and the
category $\Noetica$.

\begin{definition}[Idea Space]
\label{def:core-idea-space}
The \textbf{Idea space} $\Ideas$ is the collection of all Ideas---formal objects
representing spiritual, mental, emotional, or physical content. Formally, $\Ideas$
is realized as a graded multiset algebra:
\[
    \Ideas := \bigcup_{W \in \Worlds} \mathcal{M}(\mathcal{P}_W)
\]
where $\mathcal{M}(\mathcal{P}_W)$ denotes finite multisets over a property space
$\mathcal{P}_W$ for world $W$.
\end{definition}

\begin{definition}[Tootra-Addition Monoid]
\label{def:core-monoid}
The structure $(\Ideas, \Tadd, \mathbf{0})$ is an \textbf{abelian monoid}, where:
\begin{enumerate}[label=(\roman*)]
    \item \textbf{Binary operation} $\Tadd : \Ideas \times \Ideas \to \Ideas$ is
          multiset union:
          \[
              (X \Tadd Y)(p) := X(p) + Y(p) \quad \text{for all properties } p
          \]
    \item \textbf{Identity element} $\mathbf{0} := \varnothing$ (the empty multiset).
\end{enumerate}
This operation satisfies:
\begin{align}
    X \Tadd Y &= Y \Tadd X && \text{(Commutativity)} \\
    (X \Tadd Y) \Tadd Z &= X \Tadd (Y \Tadd Z) && \text{(Associativity)} \\
    X \Tadd \mathbf{0} &= X && \text{(Identity)}
\end{align}
\end{definition}

\begin{definition}[Category $\Noetica$]
\label{def:core-noetica}
The category $\Noetica$ is defined by:

\textbf{Objects:}
\[
    \mathrm{Ob}(\Noetica) := \Ideas \cup \Worlds \cup \Elements \cup \Foundations
                             \cup \Acquisitions
\]

\textbf{Morphisms:} The hom-sets consist of noetic transitions---primarily the
Noetic operators and world-association morphisms:
\[
    \mathrm{Hom}(\Noetica) := \Noetics \cup \{\oplus_W, \ominus_W : W \in \Worlds\}
\]

\textbf{Composition:} For morphisms $f : X \to Y$ and $g : Y \to Z$:
\[
    g \circ f : X \to Z
\]
is standard function composition (apply $f$ first, then $g$).

\textbf{Identity:} For each object $X$, the identity morphism is:
\[
    \mathrm{id}_X := \noe{0}|_X
\]
\end{definition}

\begin{proposition}[Category Axioms]
\label{prop:noetica-category}
$\Noetica$ satisfies the category axioms:
\begin{enumerate}[label=(\alph*)]
    \item \textbf{Associativity}: $(h \circ g) \circ f = h \circ (g \circ f)$,
    \item \textbf{Identity}: $\mathrm{id}_Y \circ f = f = f \circ \mathrm{id}_X$
          for all $f : X \to Y$.
\end{enumerate}
\end{proposition}

\begin{definition}[Noetic Monoid]
\label{def:core-noetic-monoid}
The structure $(\Noetics, \circ, \noe{0})$ forms a \textbf{monoid} where composition
$\circ$ is associative and $\noe{0}$ is the identity. The $10 \times 10$ composition
table is canonically specified in v6.1.
\end{definition}

%% ============================================================================
%% SECTION 3: THE RPM MONAD
%% ============================================================================
\section{The RPM Monad}
\label{sec:rpm-monad}

The Recursive Prerequisite Model (RPM) is formalized as a monad encoding prerequisite-dependent
computation with potential failure.

\begin{definition}[RPM Monad Structure]
\label{def:rpm-monad}
The \textbf{RPM monad} is a triple $(\RPM, \rpmret, \rpmbind)$ where:
\begin{enumerate}[label=(\roman*)]
    \item $\RPM : \Noetica \to \Noetica$ is an endofunctor assigning to each object $X$
          the type of RPM computations returning $X$:
          \[
              \RPM(X) := (\mathrm{AcqState} \to (X + \Failed) \times \mathrm{AcqState})
          \]
          where $\Failed$ is the failure token and $\mathrm{AcqState} = 2^{\Acquisitions}$.

    \item $\rpmret : X \to \RPM(X)$ is the \textbf{unit} (return):
          \[
              \rpmret\; v := \lambda \sigma.\, (\mathrm{inl}\; v,\, \sigma)
          \]
          It lifts a pure value into an RPM computation that succeeds immediately.

    \item $\rpmbind : \RPM(X) \times (X \to \RPM(Y)) \to \RPM(Y)$ is the \textbf{bind}:
          \[
              m \rpmbind f := \lambda \sigma.\,
              \begin{cases}
                  f(v)(\sigma') & \text{if } m(\sigma) = (\mathrm{inl}\; v, \sigma') \\
                  (\mathrm{inr}\; \Failed, \sigma') & \text{if } m(\sigma) = (\mathrm{inr}\; \Failed, \sigma')
              \end{cases}
          \]
          Sequentially composes RPM computations, propagating failure.
\end{enumerate}
\end{definition}

\begin{theorem}[RPM Monad Laws]
\label{thm:rpm-monad-laws}
The RPM monad $(\RPM, \rpmret, \rpmbind)$ satisfies the monad axioms:

\textbf{Left Unit:}
\[
    \rpmret\; x \rpmbind f \;=\; f(x)
\]

\textbf{Right Unit:}
\[
    m \rpmbind \rpmret \;=\; m
\]

\textbf{Associativity:}
\[
    (m \rpmbind f) \rpmbind g \;=\; m \rpmbind (\lambda x.\, f(x) \rpmbind g)
\]
\end{theorem}

\begin{proof}[Proof Sketch]
\textit{Left Unit}: $(\rpmret\; x \rpmbind f)(\sigma) = f(x)(\sigma)$ since
$\rpmret\; x$ returns $(\mathrm{inl}\; x, \sigma)$ immediately.

\textit{Right Unit}: $(m \rpmbind \rpmret)(\sigma)$: if $m(\sigma) = (\mathrm{inl}\; v, \sigma')$,
then $\rpmret\; v(\sigma') = (\mathrm{inl}\; v, \sigma')$, recovering $m(\sigma)$.

\textit{Associativity}: Both sides thread state identically through $m$, then $f$,
then $g$, with failure short-circuiting at each stage.
\end{proof}

%% ============================================================================
%% SECTION 4: FINITE FRACTALS
%% ============================================================================
\section{Finite Fractals}
\label{sec:finite-fractals}

Noetic fractals formalize recursive self-application of noetic transformations.

\begin{definition}[Noetic Fractal]
\label{def:noetic-fractal-core}
A \textbf{noetic fractal} $\Frac$ is a finite sequence of noetic indices:
\[
    \Frac := \langle k_1 : k_2 : \cdots : k_n \rangle
\]
where $k_i \in \{0, 1, \ldots, 9\}$ and $n \geq 1$. As an operator:
\[
    \Frac : \Ideas \to \Ideas, \qquad
    \Frac(X) := \noe{k_1}(\noe{k_2}(\cdots \noe{k_n}(X) \cdots))
\]
Application is right-to-left (innermost first).
\end{definition}

\begin{definition}[Fractal as Endofunctor]
\label{def:fractal-endofunctor}
Each fractal $\Frac$ induces an endofunctor $\Frac : \Noetica \to \Noetica$ by:
\begin{itemize}
    \item \textbf{On objects}: $\Frac(X) := \noe{k_1} \circ \cdots \circ \noe{k_n}(X)$,
    \item \textbf{On morphisms}: $\Frac(f) := \noe{k_1} \circ \cdots \circ \noe{k_n} \circ f
          \circ \noe{k_n}^{-1} \circ \cdots \circ \noe{k_1}^{-1}$ (conjugation).
\end{itemize}
The \emph{self-similarity} property is encoded by the equation:
\[
    \Frac \circ \Frac = \Frac^2
\]
which generalizes to $n$-fold iteration.
\end{definition}

\begin{definition}[Fractal Iteration]
\label{def:fractal-iteration}
For $n \in \mathbb{N}$, define the $n$-fold iteration:
\[
    \Frac^n := \underbrace{\Frac \circ \Frac \circ \cdots \circ \Frac}_{n \text{ times}}
\]
with $\Frac^0 := \mathrm{id}$ (the identity functor). Explicitly:
\[
    \Frac^n(X) = \Frac(\Frac^{n-1}(X))
\]
\end{definition}

\begin{definition}[Fractal Fixed Point]
\label{def:fractal-fixpoint-core}
A \textbf{fixed point} of fractal $\Frac$ is an object $X^* \in \mathrm{Ob}(\Noetica)$
satisfying:
\[
    \Frac(X^*) = X^*
\]
The \emph{formal fixed-point equation} for fractals is:
\[
    \Frac^* = \Frac \circ \Frac^*
\]
where $\Frac^*$ denotes the limiting fractal (when convergent).
\end{definition}

\begin{definition}[Fractal Orbit and Convergence]
\label{def:fractal-orbit}
The \textbf{orbit} of $X$ under $\Frac$ is:
\[
    \mathcal{O}_\Frac(X) := \{X,\, \Frac(X),\, \Frac^2(X),\, \Frac^3(X),\, \ldots\}
\]
The fractal \textbf{converges} on $X$ if $\lim_{n \to \infty} \Frac^n(X) = X^*$ exists
in a suitable topology on $\Ideas$.
\end{definition}

%% ============================================================================
%% SECTION 5: CORE THEOREMS
%% ============================================================================
\section{Core Theorems}
\label{sec:core-theorems}

We state the three foundational theorems of TKS without full proof; detailed proofs
appear in the main body of this manual.

\begin{theorem}[RPM Monad Laws]
\label{thm:core-rpm}
The RPM monad $(\RPM, \rpmret, \rpmbind)$ satisfies the monad axioms over the category
$\Noetica$:
\begin{enumerate}[label=(\alph*)]
    \item \textbf{Left Unit}: $\rpmret\; x \rpmbind f = f(x)$,
    \item \textbf{Right Unit}: $m \rpmbind \rpmret = m$,
    \item \textbf{Associativity}: $(m \rpmbind f) \rpmbind g = m \rpmbind (\lambda x.\, f(x) \rpmbind g)$.
\end{enumerate}
These laws ensure that prerequisite-dependent computations compose coherently, with
failure propagating correctly through chains of dependent operations.
\end{theorem}

\begin{theorem}[ACBE Functor]
\label{thm:core-acbe}
The Acquisition-Completion-Binding-Emission functor
\[
    \mathsf{ACBE} : \Noetica \longrightarrow \Acq
\]
preserves composition and identities:
\begin{enumerate}[label=(\alph*)]
    \item \textbf{Identity Preservation}: $\mathsf{ACBE}(\mathrm{id}_X) = \mathrm{id}_{\mathsf{ACBE}(X)}$,
    \item \textbf{Composition Preservation}: $\mathsf{ACBE}(g \circ f) = \mathsf{ACBE}(g) \circ \mathsf{ACBE}(f)$.
\end{enumerate}
Here $\Acq$ is the category of acquisition states and transitions. The ACBE functor
encodes the manifestation process from archetypal cause to material effect via the
composition $\noe{8} \circ (-) \circ \noe{9}$ (Above-morphism-Below).
\end{theorem}

\begin{theorem}[Fractal Fixed Point --- Kleene Convergence]
\label{thm:core-fixpoint}
Let $\Frac : \Noetica \to \Noetica$ be a finite noetic fractal satisfying the
\emph{Scott-continuity} condition (preserving directed suprema). Then the
Kleene iteration converges:
\[
    \fix(\Frac) \;=\; \bigsqcup_{n < \omega} \Frac^n(\bot)
\]
where $\bot$ is the bottom element of the noetic domain and $\bigsqcup$ denotes
the directed supremum. Explicitly:
\[
    \fix(\Frac) = \bot \sqcup \Frac(\bot) \sqcup \Frac^2(\bot) \sqcup \cdots
\]
This fixed point is the \emph{least} solution to the equation $X = \Frac(X)$ and
represents the stable attractor of iterated fractal application.
\end{theorem}

\begin{remark}[Domain-Theoretic Foundation]
\label{rem:domain-theory}
Theorem~\ref{thm:core-fixpoint} relies on the domain-theoretic structure of $\Noetica$
established in Part~VI. The noetic domain $\NoeticDomain$ forms a dcpo (directed-complete
partial order) under the information ordering $\sqsubseteq$. Scott-continuous functors
on dcpos admit least fixed points by the Kleene fixed-point theorem.
\end{remark}

\begin{corollary}[Stabilizing Fractals Converge]
\label{cor:stabilizing-convergence}
A fractal $\Frac$ is \emph{stabilizing} if and only if for all compatible inputs $X$:
\[
    \exists\, X^* \in \Ideas : \lim_{n \to \infty} \Frac^n(X) = X^*
\]
Examples include the MVR fractal $\langle 1 : 4 : 7 \rangle$ (Mind-Vibration-Rhythm)
and the Deep Mind fractal $\langle 1 : 1 : 1 \rangle$.
\end{corollary}

%% ============================================================================
%% SECTION 6: DENOTATIONAL SEMANTICS OF THE CORE CALCULUS
%% ============================================================================
\section{Denotational Semantics of the Core Calculus}
\label{sec:core-denotational}

We now establish a denotational semantics for the TKS core calculus, mapping syntactic
constructs to elements of a domain-theoretic model. This provides a compositional
interpretation independent of evaluation order.

\begin{definition}[Core Language Syntax]
\label{def:core-lang-syntax}
The \textbf{core language} $\mathcal{L}_{\mathrm{core}}$ is generated by the grammar:
\[
e ::= \iota \mid w \vdash \iota \mid \Foundations_i(\iota) \mid \rpmret(e) \mid
      e_1 \rpmbind e_2 \mid \Frac^n(e) \mid \fix(\Frac)
\]
where $\iota$ ranges over ideas, $w \in \Worlds$, $\Foundations_i \in \Foundations$,
$\Frac$ is a fractal, and $n \in \mathbb{N}$.
\end{definition}

\begin{definition}[Semantic Domain]
\label{def:semantic-domain}
The \textbf{semantic domain} $\mathbf{D}$ is a dcpo (directed-complete partial order)
constructed as follows:
\[
\mathbf{D} := \Ideas_\bot \times \mathrm{AcqState}_\bot
\]
where $\Ideas_\bot = \Ideas \cup \{\bot\}$ is the lifted Idea space, and
$\mathrm{AcqState}_\bot = 2^{\Acquisitions} \cup \{\bot\}$ is the lifted acquisition state.
The ordering is componentwise: $(X_1, \sigma_1) \sqsubseteq (X_2, \sigma_2)$ iff
$X_1 \sqsubseteq X_2$ and $\sigma_1 \sqsubseteq \sigma_2$.
\end{definition}

\begin{definition}[Semantic Valuation Function]
\label{def:sem-valuation}
The \textbf{semantic function} $\sem{-} : \mathcal{L}_{\mathrm{core}} \to \mathbf{D}$
is defined inductively:

\medskip
\textbf{(1) Ideas:}
\[
\sem{\iota} := (\iota, \sigma_0) \in \Ideas_\bot \times \mathrm{AcqState}_\bot
\]
where $\sigma_0$ is the initial acquisition state.

\medskip
\textbf{(2) World Membership:}
\[
\sem{w \vdash \iota} :=
\begin{cases}
(\iota, \sigma_0) & \text{if } \mathrm{world}(\iota) = w \\
(\bot, \sigma_0) & \text{otherwise}
\end{cases}
\]
The predicate $\mathrm{world}(\iota) = w$ holds when idea $\iota$ is associated with world $w$.

\medskip
\textbf{(3) Foundation Application:}
\[
\sem{\Foundations_i(\iota)} := (\Foundations_i \cdot \iota, \sigma_0[\mathsf{D}_i \mapsto \mathtt{true}])
\]
where $\Foundations_i \cdot \iota$ denotes the application of Foundation $\Foundations_i$
to idea $\iota$, filtering or projecting through the domain's organizational principle.

\medskip
\textbf{(4) RPM Return:}
\[
\sem{\rpmret(e)} := \lambda \sigma.\, (\mathrm{inl}\,\sem{e}_{\mathrm{val}}, \sigma)
\]
where $\sem{e}_{\mathrm{val}}$ extracts the value component from $\sem{e}$.

\medskip
\textbf{(5) RPM Bind (Sequencing):}
\[
\sem{m \rpmbind f} := \lambda \sigma.\,
\begin{cases}
\sem{f}(v)(\sigma') & \text{if } \sem{m}(\sigma) = (\mathrm{inl}\, v, \sigma') \\
(\mathrm{inr}\,\Failed, \sigma') & \text{if } \sem{m}(\sigma) = (\mathrm{inr}\,\Failed, \sigma')
\end{cases}
\]
This propagates state through sequential composition and short-circuits on failure.

\medskip
\textbf{(6) Fractal Iteration:}
\[
\sem{\Frac^n(\iota)} := \underbrace{\sem{\Frac} \circ \sem{\Frac} \circ \cdots \circ \sem{\Frac}}_{n \text{ times}}(\sem{\iota})
\]
where $\sem{\Frac} = \sem{\noe{k_1}} \circ \cdots \circ \sem{\noe{k_m}}$ for
$\Frac = \langle k_1 : \cdots : k_m \rangle$.

\medskip
\textbf{(7) Fractal Fixed Point:}
\[
\sem{\fix(\Frac)} := \bigsqcup_{n < \omega} \sem{\Frac}^n(\bot)
\]
the least fixed point computed via Kleene iteration.
\end{definition}

\begin{proposition}[Compositionality]
\label{prop:denotational-compositionality}
The semantic function $\sem{-}$ is compositional: the meaning of a compound expression
depends only on the meanings of its immediate subexpressions.
\end{proposition}

\begin{theorem}[Semantic Domain Properties]
\label{thm:sem-domain-props}
The semantic domain $\mathbf{D}$ satisfies:
\begin{enumerate}[label=(\roman*)]
    \item $\mathbf{D}$ is a pointed dcpo with bottom element $(\bot, \bot)$.
    \item All semantic operators $\sem{\Frac}$, $\sem{\noe{k}}$ are Scott-continuous.
    \item The RPM operations preserve continuity: if $f$ is continuous, so is
          $\sem{m \rpmbind f}$ as a function of $\sem{m}$.
\end{enumerate}
\end{theorem}

%% ============================================================================
%% SECTION 7: OPERATIONAL SEMANTICS OF THE CORE CALCULUS
%% ============================================================================
\section{Operational Semantics of the Core Calculus}
\label{sec:core-operational}

We define a big-step operational semantics that specifies how core calculus expressions
evaluate to values. This provides an executable specification complementing the
denotational account.

\begin{definition}[Values and Configurations]
\label{def:operational-values}
\textbf{Values} are terminal expressions:
\[
v ::= \iota \mid \rpmret(v) \mid \mathrm{inr}\,\Failed
\]
A \textbf{configuration} is a pair $(e, \sigma)$ of an expression $e$ and an acquisition
state $\sigma \in 2^{\Acquisitions}$. We write $e \bigstep v$ for pure evaluation and
$(e, \sigma) \bigstep (v, \sigma')$ for stateful evaluation.
\end{definition}

\begin{definition}[Big-Step Evaluation Rules]
\label{def:bigstep-rules}
The big-step relation $\bigstep$ is defined by the following inference rules:

\medskip
\textbf{Rule (Val) --- Values:}
\[
\infer[\textsc{Val}]{v \bigstep v}{}
\]
Values evaluate to themselves.

\medskip
\textbf{Rule (World) --- World Membership:}
\[
\infer[\textsc{World-Ok}]{(w \vdash \iota, \sigma) \bigstep (\iota, \sigma)}{\mathrm{world}(\iota) = w}
\]
\[
\infer[\textsc{World-Fail}]{(w \vdash \iota, \sigma) \bigstep (\bot, \sigma)}{\mathrm{world}(\iota) \neq w}
\]

\medskip
\textbf{Rule (Found) --- Foundation Application:}
\[
\infer[\textsc{Found}]{(\Foundations_i(\iota), \sigma) \bigstep (\Foundations_i \cdot \iota,\, \sigma[\mathsf{D}_i \mapsto \mathtt{true}])}{}
\]

\medskip
\textbf{Rule (Ret) --- RPM Return:}
\[
\infer[\textsc{Ret}]{(\rpmret(v), \sigma) \bigstep (\mathrm{inl}\, v, \sigma)}{}
\]
The unit lifts a value into a successful RPM computation.

\medskip
\textbf{Rule (Bind-Ok) --- RPM Bind (Success):}
\[
\infer[\textsc{Bind-Ok}]{(m \rpmbind f, \sigma) \bigstep (v', \sigma'')}{%
    (m, \sigma) \bigstep (\mathrm{inl}\, v, \sigma') & (f(v), \sigma') \bigstep (v', \sigma'')}
\]
If the first computation succeeds with value $v$ and updated state $\sigma'$,
then apply the continuation $f$ to $v$ in state $\sigma'$.

\medskip
\textbf{Rule (Bind-Fail) --- RPM Bind (Failure Propagation):}
\[
\infer[\textsc{Bind-Fail}]{(m \rpmbind f, \sigma) \bigstep (\mathrm{inr}\,\Failed, \sigma')}{%
    (m, \sigma) \bigstep (\mathrm{inr}\,\Failed, \sigma')}
\]
If the first computation fails, the failure propagates without invoking the continuation.

\medskip
\textbf{Rule (Frac-Zero) --- Fractal Base Case:}
\[
\infer[\textsc{Frac-0}]{\Frac^0(\iota) \bigstep \iota}{}
\]
Zero iterations yield the input unchanged.

\medskip
\textbf{Rule (Frac-Step) --- Fractal Iteration Step:}
\[
\infer[\textsc{Frac-Step}]{\Frac^{n+1}(\iota) \bigstep v'}{%
    \Frac^n(\iota) \bigstep v & \Frac(v) \bigstep v'}
\]
Iteration proceeds by evaluating $n$ steps, then applying $\Frac$ once more.

\medskip
\textbf{Rule (Noetic) --- Noetic Application:}
\[
\infer[\textsc{Noetic}]{\noe{k}(\iota) \bigstep \noe{k} \cdot \iota}{}
\]
where $\noe{k} \cdot \iota$ denotes the result of applying noetic $\noe{k}$ to idea $\iota$.

\medskip
\textbf{Rule (Fix) --- Fixed Point Unfolding:}
\[
\infer[\textsc{Fix}]{\fix(\Frac) \bigstep v}{\Frac(\fix(\Frac)) \bigstep v}
\]
The fixed point satisfies $\fix(\Frac) = \Frac(\fix(\Frac))$.
\end{definition}

\begin{proposition}[Determinism]
\label{prop:bigstep-determinism}
The big-step relation is deterministic: if $e \bigstep v_1$ and $e \bigstep v_2$,
then $v_1 = v_2$.
\end{proposition}

\begin{proof}[Proof Sketch]
By induction on the derivation of $e \bigstep v_1$. Each rule has disjoint applicability
conditions (e.g., \textsc{Bind-Ok} vs. \textsc{Bind-Fail} are distinguished by the
outcome of the first premise), ensuring at most one rule applies to any configuration.
\end{proof}

%% ============================================================================
%% SECTION 8: SOUNDNESS AND ADEQUACY
%% ============================================================================
\section{Soundness and Adequacy}
\label{sec:core-soundness}

We establish that the denotational and operational semantics agree, ensuring that
the domain-theoretic model faithfully represents the computational behavior of
the core calculus.

\begin{theorem}[Soundness]
\label{thm:core-soundness}
If $e \bigstep v$, then $\sem{e} = \sem{v}$.

\noindent
Equivalently, operational evaluation preserves denotational meaning: a program's
mathematical semantics equals the semantics of its computed result.
\end{theorem}

\begin{proof}[Proof Sketch]
By induction on the derivation $e \bigstep v$.

\textit{Case} \textsc{Val}: Immediate, since $\sem{v} = \sem{v}$.

\textit{Case} \textsc{Ret}: We have $\rpmret(v) \bigstep \mathrm{inl}\, v$.
By definition, $\sem{\rpmret(v)} = \lambda \sigma.\, (\mathrm{inl}\,\sem{v}, \sigma)$,
which equals $\sem{\mathrm{inl}\, v}$ as a semantic object.

\textit{Case} \textsc{Bind-Ok}: Suppose $(m, \sigma) \bigstep (\mathrm{inl}\, v, \sigma')$
and $(f(v), \sigma') \bigstep (v', \sigma'')$. By IH, $\sem{m}(\sigma) = \sem{\mathrm{inl}\, v}$
at state $\sigma'$, and $\sem{f(v)}(\sigma') = \sem{v'}$ at state $\sigma''$.
The definition of $\sem{m \rpmbind f}$ then yields $\sem{v'}$ at $\sigma''$.

\textit{Case} \textsc{Bind-Fail}: If $(m, \sigma) \bigstep (\mathrm{inr}\,\Failed, \sigma')$,
by IH $\sem{m}(\sigma)$ is the failure case. The semantics of bind propagates this failure.

\textit{Case} \textsc{Frac-Step}: By IH on both premises and compositionality of $\sem{-}$.

\textit{Case} \textsc{Fix}: The fixed-point equation $\fix(\Frac) = \Frac(\fix(\Frac))$
holds both operationally (by the rule) and denotationally (by Kleene's theorem).
\end{proof}

\begin{theorem}[Adequacy]
\label{thm:core-adequacy}
If $\sem{e} \neq \bot$, then there exists a value $v$ such that $e \bigstep v$.

\noindent
Equivalently, if a program has a defined (non-bottom) denotation, then operational
evaluation terminates with a result.
\end{theorem}

\begin{proof}[Proof Sketch]
By induction on the structure of expressions in $\mathcal{L}_{\mathrm{core}}$.

\textit{Case} $e = \iota$ (Idea): $\sem{\iota} = (\iota, \sigma_0) \neq \bot$,
and $\iota \bigstep \iota$ by \textsc{Val}.

\textit{Case} $e = w \vdash \iota$: If $\sem{w \vdash \iota} \neq \bot$, then
$\mathrm{world}(\iota) = w$, so \textsc{World-Ok} applies.

\textit{Case} $e = m \rpmbind f$: If $\sem{m \rpmbind f}(\sigma) \neq \bot$, then either:
\begin{itemize}
    \item $\sem{m}(\sigma) = (\mathrm{inl}\, v, \sigma')$ with $\sem{f(v)}(\sigma') \neq \bot$:
          by IH on $m$ and $f(v)$, both evaluate, so \textsc{Bind-Ok} applies.
    \item $\sem{m}(\sigma) = (\mathrm{inr}\,\Failed, \sigma')$: by IH on $m$,
          $m$ evaluates to failure, so \textsc{Bind-Fail} applies.
\end{itemize}

\textit{Case} $e = \Frac^n(\iota)$: By induction on $n$. Base case $n = 0$ is immediate.
For $n + 1$, by IH $\Frac^n(\iota) \bigstep v$ for some $v$, and $\Frac(v)$ is
well-defined by the finitary nature of fractals.

\textit{Case} $e = \fix(\Frac)$: The Kleene fixed-point theorem guarantees that
$\bigsqcup_n \Frac^n(\bot)$ is reached in finite iterations for Scott-continuous $\Frac$
over the finite noetic domain. Hence operational unfolding terminates.
\end{proof}

\begin{corollary}[Semantic Equivalence]
\label{cor:semantic-equivalence}
The denotational and operational semantics define the same partial function:
\[
e \bigstep v \quad \Longleftrightarrow \quad \sem{e} = \sem{v} \neq \bot
\]
\end{corollary}

\begin{remark}[Adequacy for the Finite Core]
\label{rem:adequacy-finite}
Adequacy holds straightforwardly for the \emph{finite} core calculus because:
\begin{enumerate}[label=(\roman*)]
    \item The Idea space $\Ideas$ is finite (40 elements).
    \item All noetics and fractals operate over this finite domain.
    \item Fractal iteration either stabilizes or cycles within finite steps.
    \item RPM computations are parametrized by finite acquisition states ($2^{22}$ possibilities).
\end{enumerate}
Extension to transfinite fractals requires additional continuity conditions, established
in the full domain-theoretic treatment of Part~VI.
\end{remark}

\begin{theorem}[Type Soundness (Core Fragment)]
\label{thm:type-soundness-core}
If $\Gamma \vdash e : \tau$ is derivable in the TKS type system, then either:
\begin{enumerate}[label=(\roman*)]
    \item $e$ is a value, or
    \item $e \bigstep v$ for some value $v$ with $\Gamma \vdash v : \tau$.
\end{enumerate}
Well-typed programs do not get stuck.
\end{theorem}

\begin{proof}[Proof Sketch]
By the standard technique of progress and preservation:
\begin{itemize}
    \item \textbf{Progress}: A well-typed, non-value expression can always take an
          evaluation step (some rule in Definition~\ref{def:bigstep-rules} applies).
    \item \textbf{Preservation}: Evaluation preserves typing: if $\Gamma \vdash e : \tau$
          and $e \bigstep v$, then $\Gamma \vdash v : \tau$.
\end{itemize}
The finite core fragment's simple type structure (Ideas, RPM-wrapped values, fractals)
ensures both properties hold.
\end{proof}

%% ============================================================================
%% CLOSING SUMMARY
%% ============================================================================

\begin{tcolorbox}[colback=gray!5!white,colframe=gray!75!black,title=\textbf{Summary: The TKS Core Calculus},breakable]
The minimal TKS formal system comprises:
\begin{enumerate}
    \item \textbf{Finite ontology}: 4 Worlds, 40 Elements, 10 Noetics, 7 Foundations,
          28 Sub-Foundations, 22 Acquisitions.
    \item \textbf{Algebraic structure}: The abelian monoid $(\Ideas, \Tadd, \mathbf{0})$
          and the category $\Noetica$ of noetic objects and transitions.
    \item \textbf{Computational monad}: The RPM monad $(\RPM, \rpmret, \rpmbind)$ encoding
          prerequisite-dependent computation with failure.
    \item \textbf{Recursive dynamics}: Noetic fractals as endofunctors with self-similarity,
          admitting fixed points via Kleene iteration.
    \item \textbf{Manifestation functor}: ACBE mapping archetypal causes to manifest effects.
    \item \textbf{Denotational semantics}: Compositional mapping $\sem{-} : \mathcal{L}_{\mathrm{core}} \to \mathbf{D}$
          to a dcpo semantic domain.
    \item \textbf{Operational semantics}: Big-step evaluation relation $\bigstep$ specifying
          how expressions reduce to values.
    \item \textbf{Soundness and adequacy}: The operational and denotational semantics agree---evaluation
          preserves meaning, and defined programs terminate.
\end{enumerate}
This core suffices to formalize consciousness transformations, prerequisite reasoning,
and iterative refinement processes within the TKS framework, with rigorous semantic foundations.
\end{tcolorbox}

